
\documentclass{article}
\usepackage{colortbl}
\usepackage{makecell}
\usepackage{multirow}
\usepackage{supertabular}

\begin{document}

\newcounter{utterance}

\centering \large Interaction Transcript for game `cladder', experiment `full\_v1.5\_default', episode 8751 with selfrefine{-}Qwen3.5{-}27B.
\vspace{24pt}

{ \footnotesize  \setcounter{utterance}{1}
\setlength{\tabcolsep}{0pt}
\begin{supertabular}{c@{$\;$}|p{.15\linewidth}@{}p{.15\linewidth}p{.15\linewidth}p{.15\linewidth}p{.15\linewidth}p{.15\linewidth}}
    \# & \multicolumn{2}{c}{Player} && \multicolumn{2}{c}{Game Master} \\
    \hline

    \theutterance \stepcounter{utterance}  
    & & & \multicolumn{4}{p{0.6\linewidth}}{
        \cellcolor[rgb]{0.9,0.9,0.9}{
            \makecell[{{p{\linewidth}}}]{
                \texttt{\tiny{[P1$\langle$GM]}}
                \texttt{You are an expert in causal inference. The following question is not a typical commonsense query, but rather a meticulously designed question created by a professor specializing in causal inference, intended to assess the students' mastery of the course content.} \\
\texttt{Imagine a self{-}contained, hypothetical world with only the following conditions, and without any unmentioned factors or causal relationships: Yupt has a direct effect on zupj. Zupj has a direct effect on muvq. For those who are not yupt, the probability of muvq is 25\%. For those who are yupt, the probability of muvq is 25\%. For those who are yupt, would it be more likely to see muvq if the individual was not yupt?} \\
\texttt{Start your answer with 'Yes' or 'No', followed by additional reasoning or evidence to support your explanation.} \\
            }
        }
    }
     \\ \\

    \theutterance \stepcounter{utterance}  
    & & \multicolumn{4}{p{0.6\linewidth}}{
        \cellcolor[rgb]{0.95,0.95,0.95}{
            \makecell[{{p{\linewidth}}}]{
                \texttt{\tiny{[P1$\leftrightarrow$P]}}
                \texttt{{[}DRAFT{]}} \\
\texttt{No} \\
\\ 
\texttt{The reasoning is as follows: The question asks whether an individual who is currently "yupt" would be more likely to see "muvq" if they were instead "not yupt." This is a counterfactual query comparing the probability of the outcome under two different states of the treatment variable (Yupt).} \\
\\ 
\texttt{According to the conditions provided:} \\
\texttt{1.  The probability of `muvq` given `yupt` is 25\% (\$P(\textbackslash{}text\{muvq\} | \textbackslash{}text\{yupt\}) = 0.25\$).} \\
\texttt{2.  The probability of `muvq` given `not yupt` is 25\% (\$P(\textbackslash{}text\{muvq\} | \textbackslash{}neg\textbackslash{}text\{yupt\}) = 0.25\$).} \\
\\ 
\texttt{Although the problem states there is a causal chain (\$Yupt \textbackslash{}to Zupj \textbackslash{}to Muvq\$), the observed marginal probabilities indicate that the net effect of `yupt` on `muvq` is zero. In a self{-}contained world with no unmentioned factors, if the probability of the outcome is identical regardless of the treatment status, the treatment has no causal effect on the outcome (the Average Treatment Effect is 0).} \\
\\ 
\texttt{Therefore, if an individual who is `yupt` were hypothetically changed to `not yupt`, their probability of experiencing `muvq` would remain 25\%. It would not be *more* likely; it would be exactly as likely.} \\
            }
        }
    }
    & \\ \\

    \theutterance \stepcounter{utterance}  
    & & \multicolumn{4}{p{0.6\linewidth}}{
        \cellcolor[rgb]{0.95,0.95,0.95}{
            \makecell[{{p{\linewidth}}}]{
                \texttt{\tiny{[P1$\leftrightarrow$P]}}
                \texttt{{[}SELF{-}FEEDBACK{]}} \\
\texttt{FORMAT: problem} \\
\texttt{RULES: ok} \\
\texttt{STRATEGY: ok} \\
\texttt{REASON: The answer includes explanatory text before the required starting word "No", violating the instruction to start the answer immediately with 'Yes' or 'No'.} \\
            }
        }
    }
    & \\ \\

    \theutterance \stepcounter{utterance}  
    & \multicolumn{4}{p{0.6\linewidth}}{
        \cellcolor[rgb]{0.561,0.737,0.561}{
            \makecell[{{p{\linewidth}}}]{
                \texttt{\tiny{[P1$\rangle$GM]}}
                \texttt{No} \\
\texttt{The question asks whether an individual who is currently "yupt" would be more likely to see "muvq" if they were instead "not yupt." This is a counterfactual query comparing the probability of the outcome under two different states of the treatment variable (Yupt).} \\
\\ 
\texttt{According to the conditions provided:} \\
\texttt{1. The probability of `muvq` given `yupt` is 25\% (\$P(\textbackslash{}text\{muvq\} | \textbackslash{}text\{yupt\}) = 0.25\$).} \\
\texttt{2. The probability of `muvq` given `not yupt` is 25\% (\$P(\textbackslash{}text\{muvq\} | \textbackslash{}neg\textbackslash{}text\{yupt\}) = 0.25\$).} \\
\\ 
\texttt{Although the problem states there is a causal chain (\$Yupt \textbackslash{}to Zupj \textbackslash{}to Muvq\$), the observed marginal probabilities indicate that the net effect of `yupt` on `muvq` is zero. In a self{-}contained world with no unmentioned factors, if the probability of the outcome is identical regardless of the treatment status, the treatment has no causal effect on the outcome (the Average Treatment Effect is 0).} \\
\\ 
\texttt{Therefore, if an individual who is `yupt` were hypothetically changed to `not yupt`, their probability of experiencing `muvq` would remain 25\%. It would not be *more* likely; it would be exactly as likely.} \\
            }
        }
    }
    & & \\ \\

    \theutterance \stepcounter{utterance}  
    & & & \multicolumn{2}{p{0.3\linewidth}}{
        \cellcolor[rgb]{0.9,0.9,0.9}{
            \makecell[{{p{\linewidth}}}]{
                \texttt{\tiny{[GM$|$GM]}}
                \texttt{No} \\
            }
        }
    }
    & & \\ \\

    \theutterance \stepcounter{utterance}  
    & & & \multicolumn{2}{p{0.3\linewidth}}{
        \cellcolor[rgb]{0.9,0.9,0.9}{
            \makecell[{{p{\linewidth}}}]{
                \texttt{\tiny{[GM$|$GM]}}
                \texttt{no} \\
            }
        }
    }
    & & \\ \\

    \theutterance \stepcounter{utterance}  
    & & & \multicolumn{2}{p{0.3\linewidth}}{
        \cellcolor[rgb]{0.9,0.9,0.9}{
            \makecell[{{p{\linewidth}}}]{
                \texttt{\tiny{[GM$|$GM]}}
                \texttt{game\_result = WIN} \\
            }
        }
    }
    & & \\ \\

\end{supertabular}
}

\end{document}
